\documentclass[t]{beamer}
\usepackage[utf8]{inputenc}
\usepackage[francais]{babel}
\usepackage{tikz,pgf}
\usepackage{wrapfig}
\usepackage{times}
\usepackage[absolute,overlay]{textpos}
\usepackage{url}
% \usepackage[texcoord,grid,gridunit=mm,gridcolor=red!10,subgridcolor=green!10]{eso-pic}



% Display code
\usepackage{listings}
    \definecolor{dkgreen}{rgb}{0,0.6,0}
    \definecolor{gray}{rgb}{0.5,0.5,0.5}
    \definecolor{mallow}{rgb}{0.58,0,0.82}

    \lstset{
        language=C++,
        aboveskip=3mm,
        belowskip=3mm,
        showstringspaces=false,
        columns=flexible,
        basicstyle={\small\ttfamily},
        numbers=none,
        numberstyle=\tiny\color{gray},
        keywordstyle=\color{blue},
        commentstyle=\color{dkgreen},
        stringstyle=\color{mallow},
        breaklines=true,
        breakatwhitespace=true,
        tabsize=4,
        frame=single,
        literate=%
            {é}{{\'e}}1
            {É}{{\'E}}1
            {è}{{\`e}}1
            {à}{{\`a}}1
            {È}{{\`E}}1
            {À}{{\`A}}1
            {ê}{{\^e}}1
            {â}{{\^a}}1
            {î}{{\^\i}}1
            {ô}{{\^o}}1
            {û}{{\^u}}1
            {Ê}{{\^E}}1
            {Â}{{\^A}}1
            {Î}{{\^I}}1
            {Ô}{{\^O}}1
            {Û}{{\^U}}1
            {ë}{{\"e}}1
            {ï}{{\"\i}}1
            {ü}{{\"u}}1
            {Ë}{{\"E}}1
            {Ï}{{\"I}}1
            {Ü}{{\"U}}1
            {ç}{{\c c}}1
            {Ç}{{\c C}}1
            {æ}{{\ae}}1
            {Æ}{{\AE}}1
            {œ}{{\oe}}1
            {Œ}{{\OE}}1
            {Å}{{\AA}}1
    }




\definecolor{bleu}{rgb}{0.282, 0.439, 0.643}

\usecolortheme[named=bleu]{structure}
\useinnertheme[shadow]{rounded}

\hypersetup{pdfpagemode=FullScreen}

\setbeamersize{text margin left=0.6cm}
\setbeamersize{text margin right=0.4cm}
\setbeamersize{text margin left=.5cm, text margin right=.5cm}
\setbeamertemplate{navigation symbols}{}

\AtBeginSection[]{
    \begin{frame}
  	 \frametitle{\vspace*{-1cm} \begin{flushright} \ \ \vskip-0.25cm \end{flushright} }
        \vspace*{-0.75cm}
	    \begin{center}
            \centering
            \tableofcontents[currentsection]
	    \end{center}
    \end{frame}
}

\begin{document}
    \begin{frame}
        \pgfputat{\pgfxy(-0.505,0.685)}{\pgfbox[left,top]{\pgfimage[width=\paperwidth,height=\paperheight ]{./assets/UTT_garde_empty.pdf}}}
        \vspace*{2.5cm}

        \begin{center}
            \textcolor{white}{
                {\large Soutenance de TX} \\
                {\footnotesize   Printemps 2015}\\
            \vspace*{0.5cm}
                {\footnotesize \textbf{Développement d'une librairie de fonctions de chiffrement en C}}\\
            \vspace*{0.25cm}
                {\footnotesize Pergoud Florent\\ Labate Aurélien}\\
          }
        \end{center}
    \end{frame}

\setbeamertemplate{background canvas}{
    \includegraphics[width=\paperwidth,height=\paperheight]{./assets/UTT_fond.pdf}}
    \addtobeamertemplate{footline}{
    \vspace{-0.5cm}
    \begin{flushright}
    \insertframenumber/\inserttotalframenumber
    \end{flushright}
}


\section{Introduction}
    \begin{frame}[c]
        \begin{center}
            \frametitle{\vspace*{-1.0cm} \begin{flushright} Introduction \ \ \vskip-0.25cm \end{flushright} }
            \vspace*{-0.75cm}
            \begin{itemize}
                \item Acquérir des connaissances en cryptographie 
                \item Création d'une librairie en C avec des fonctions de cryptographie
                \item Faire un programme d'exemple avec un menu pour présenter les algorithmes
            \end{itemize}
        \end{center}
    \end{frame}


\section{Environnement de programmation}


\defverbatim[colored]\makeset{\begin{lstlisting}[language=C++]
    mpz_t a, b, c, res;
    mpz_init_set_ui(a, 42);
    mpz_init_set_str(b, "3172371623812", 10);
    mpz_init_set_str(c, "1A", 16);

    mpz_mul(res, b, c);
    mpz_add(res, a, res);

    printf("%s\n", mpz_get_str (NULL, 10, res));
\end{lstlisting}}

    \begin{frame}[c]
        \frametitle{\vspace*{-1.0cm} \begin{flushright}Langage de programmation \ \ \vskip-0.25cm \end{flushright} }
        \begin{itemize}
            \item Projet originalement proposé en C avec GMP
            \item Mais la syntaxe est très lourde en utilisant la librairie GMP en C
            \item Par exemple pour $42+ 3172371623812\times 0x1A$ :
        \end{itemize}
        \makeset
    \end{frame}


\defverbatim[colored]\makeset{\begin{lstlisting}[language=C++]
    mpz_class a, b, c, res;
    a = 42;
    b = "3172371623812";
    c = "0x1A";

    res = a + b * c;

    std::cout << "Résultat : " << res << std::endl;
\end{lstlisting}}


    \begin{frame}[c]
        \frametitle{\vspace*{-1.0cm} \begin{flushright}Langage de programmation \ \ \vskip-0.25cm \end{flushright} }
        \begin{itemize}
            \item Alors qu'avec la version C++ de GMP :
            \makeset
            \item Cependant, le C++ est moins connu est peut gêner
        \end{itemize}
    \end{frame}


\defverbatim[colored]\makeset{\begin{lstlisting}[language=C++]
    mpz_class a, res;
    a = 2;

    res = mpz_pow_ui(a.get_mpz_t(), 3);

    std::cout << "Résultat : " << res << std::endl;
\end{lstlisting}}


    \begin{frame}[c]
        \frametitle{\vspace*{-1.0cm} \begin{flushright}Bibliotheque \ \ \vskip-0.25cm \end{flushright} }
        \begin{itemize}
            \item La version C++ de GMP manque certaines fonctionnalités
            \item Ainsi pour calculer $2^3$ :

            \makeset
            \item On doit utiliser les fonction de la version C
        \end{itemize}
    \end{frame}



\defverbatim[colored]\makeset{\begin{lstlisting}[language=C++]
    mpz2_class a, res;
    a = 2;

    res = a.pow(3);

    std::cout << "Résultat : " << res << std::endl;
\end{lstlisting}}


    \begin{frame}[c]
        \frametitle{\vspace*{-1.0cm} \begin{flushright}Bibliotheque \ \ \vskip-0.25cm \end{flushright} }
        \begin{itemize}
            \item Nous avons donc créer une librairie qui évite cela tout en gardant les avantages :

            \makeset
        \end{itemize}
    \end{frame}

\section{Les algorithmes principaux}
        \begin{frame}[c]
            \frametitle{\vspace*{-1.0cm} \begin{flushright} Les algorithmes principaux \ \ \vskip-0.25cm \end{flushright} }
            \vspace*{0.5cm}
			\begin{center}
			\textbf {Les différents algorithmes}
			\end{center}
            \begin{center}
                \begin{itemize}
                    \item Exponentiation rapide
                    \item Test de primalité de Rabin-Miller
                    \item Diffie-Helmann
                    \item ElGamal
                    \item RSA
                \end{itemize}
            \end{center}
        \end{frame}

        \begin{frame}[c]
            \frametitle{\vspace*{-1.0cm} \begin{flushright} Les algorithmes principaux \ \ \vskip-0.25cm \end{flushright} }
            \vspace*{0.5cm}
             \begin{center}
             \textbf{Test de primalité de Rabin-Miller}
             \end{center}
            \begin{center}
                \begin{itemize}
                    \item Algorithme probabiliste
                    \item En entrée un entier n à tester et un entier k (précision)
                    \item En sortie "probablament premier" ou "composé"
                    \item Probabilité de se tromper :  $4^{-k}$ 
                \end{itemize}
            \end{center}
        \end{frame}

        \begin{frame}[c]
            \frametitle{\vspace*{-1.0cm} \begin{flushright} Les algorithmes principaux \ \ \vskip-0.25cm \end{flushright} }
            \vspace*{0.5cm}
             \begin{center}
             \textbf{ElGamal}
             \end{center}
            \begin{center}
                \begin{itemize}
                    \item Algorithme servant à chiffrer et déchiffrer un message 
                    \item Clé publique d'Alice : (p,g,A)
                    \item Clés privées a et b avec $A=g^a$ et $B=g^b$
                    \item Message chiffré de Bob : (M,B)
                    \item Déchiffrement $m=M/B^a$
                \end{itemize}
            \end{center}
        \end{frame}

        \begin{frame}[c]
            \frametitle{\vspace*{-1.0cm} \begin{flushright} Les algorithmes principaux \ \ \vskip-0.25cm \end{flushright} }
            \vspace*{0.5cm}
             \begin{center}
             \textbf{RSA}
             \end{center}
            \begin{center}
                \begin{itemize}
 				  \item Algorithme servant à chiffrer et déchiffrer un message 
                    \item Génération d'un couple de clés : clé publique (n,e), clé privée (n,d)
                    \item Généralement on prend un e constant
                    \item $n=p\times q$ ou p et q sont premiers
                    \item $d=e^{-1} \mod\phi(n)$ où $\phi(n) = (p-1)(q-1)$
                    \item chiffrement : $C=M^e \mod(n)$
                    \item déchiffrement : $M=C^d \mod(n)$
                \end{itemize}
            \end{center}
        \end{frame}

\section{Conclusion}
    \begin{frame}[c]
        \frametitle{\vspace*{-1.0cm} \begin{flushright} Conclusion \ \ \vskip-0.25cm \end{flushright} }
        \vspace*{-0.75cm}
        \begin{itemize}
            \item Beaucoup de connaissances acquises ou approfondies:
            \begin{itemize}
            		\item Notions et principes de crytographie
            		\item C++
            		\item GMP
            		\item \LaTeX
            		\item Gestion de projet de recherche
			\end{itemize}
             
            \item La bibliothèque a été crée et est fonctionnelle
            \item Beaucoup d'exemples pour utiliser la bibliothèque et comprendre les algorithmes de cryptographie en jeu
            
        \end{itemize}

    \end{frame}



\end{document}